\section{Upper quotas with a growing row cap in very tall arrays}\label{sec:upper-quota}

Let $d\to\infty$, $n\ge C_*d^4\log(d+1)$ for a sufficiently large absolute constant $C_*$ and $1\le r<d$, allowing $r$ to vary. Write $T=nr/d$, $D=\lfloor T\rfloor$, $H(w)=\sum_{s=0}^r\binom ds w^s$, and $A_i=[w^{di}]H(w)^n$. Let $a_i$ count binary arrays with all column sums $i$ and all row sums at most $r$. Uniformly for $0\le D-i\le Cd$ at any fixed $C$, we have

\begin{equation*}
a_i=A_i\sqrt d\,(2\pi\Lambda_*(i))^{-(d-1)/2}(1+o(1)),
\qquad
\Lambda_*(i)=\frac{i(d-r)}{d-1}.
\end{equation*}

At internal quotas set $\Gamma_k=a_k^2/(a_{k-1}a_{k+1})$ and $t_k=(a_{k-1}/a_{k+1})^{1/2}$. If $(D-k)/d\to c\in(0,\infty)$, then $\log\Gamma_{k+j}\to1/c$ for any fixed integer $j$ and $\log t_k=d\log w_k+o(1)$, where $w_k$ gives the size law of mean $dk/n$. If $(D-k)/d\to0$, then $\Gamma_k\to\infty$.

\subsection{Fourier balance uniform over upper profiles}

Fix a profile $(s_1,\ldots,s_n)$ with $0\le s_j\le r$ and total $di$, and set $e=nr-di=O(d^2)$. At least $n-e$ rows have size $r$. For each size $s$, let $f_s(\theta)=\mathbb E e^{i\sum_{a\in U}\theta_a}$ for a uniformly chosen $s$-subset $U$ of the columns. The balance probability is the pinned $(d-1)$-torus integral of $\prod_jf_{s_j}(\theta)e^{-\mathrm{i}\,i\sum_a\theta_a}$, with $\theta_d=0$ and normalized Haar measure. Common shifts leave the integrand unchanged because the prescribed total is $di$.

Set $V=\mathbf1^\perp\subset\mathbb R^d$ and $R=d^{-1}$. Project onto $V$ along the common angle. The Euclidean ball $\|y\|\le R$ in $V$ is injective on the resulting quotient: its pinned inverse is $\theta_a=y_a-y_d$ modulo $2\pi$, with all lifts of magnitude at most $2R$. The Jacobian from pinned coordinates to Euclidean hyperplane volume is $1/\sqrt d$, since the projection of their basis has Gram matrix $I_{d-1}-\mathbf1\mathbf1^\top/d$. Normalized Haar measure on the ball is therefore $\sqrt d\,(2\pi)^{-(d-1)}dy$.

Put $z_j=e^{\mathrm{i}\theta_j}$ and $D_z=\sum_j|z_j-\bar z|^2$. Generate a uniform $s$-subset by placing a uniform permutation into slots, tossing a fair coin to select one member from each of the first $m=\min(s,d-s)$ disjoint pairs, and taking $s-m$ fixed unused slots. Conditional Fourier modulus is at most $e^{-X/8}$, where $X$ sums the squared chords of the pairs, since each fair choice has modulus $\sqrt{1-|z_i-z_j|^2/4}$. A random distinct pair has mean square chord $2D_z/(d-1)$, hence $\mathbb E X=2mD_z/(d-1)$. Each chord square is at most four. For two disjoint ordered pairs the expected product is at most $d(d-1)/[(d-2)(d-3)]$ times the square of the one-pair mean, by allowing intersections in the numerator. Thus $\mathbb E X^2\le C(\mathbb E X+(\mathbb E X)^2)$. The inequality $1-e^{-X/8}\ge cX/(1+X)$ and Cauchy--Schwarz give a modulus loss at least $c\min(mD_z/d,1)$, hence at least $ca_s\min(D_z,1)$ with $a_s=s(d-s)/[d(d-1)]$.

Outside the projected ball $D_z\ge c_0R^2$. Indeed anchor at an index minimizing $|z_j-\bar z|$. The sum of squared chords to that anchor is at most $2D_z$, so the phases relative to it in $[-\pi,\pi]$ have squared Euclidean norm at most $C D_z$ by the chord-angle inequality. Removing their average can only decrease that norm. If $D_z<c_0R^2$ for small enough $c_0$, the centered representative lies inside the ball. Applying the modulus bound to the $n-e$ size-$r$ rows, with $a_r\ge1/d$, bounds the exterior integral by $\exp(-c'n/d^3)$.

On the ball let $L=\sum_{a\in U}y_a$ for a row of size $s$. By exchangeability $\mathbb E L=0$ and $\mathbb E L^2=a_s\|y\|^2$. For the fourth moment replace the subset by its complement if $s>d/2$, which negates $L$. Set $p_j=(m)_j/(d)_j$ and $Q_j=\sum_a y_a^j$. Expansion using $Q_1=0$ gives

\begin{equation*}
\mathbb E L^4=(p_1-7p_2+12p_3-6p_4)Q_4+(3p_2-6p_3+3p_4)Q_2^2.
\end{equation*}

The index patterns $4,3+1,2+2,2+1+1$ give $p_1Q_4$, $-4p_2Q_4$, $3p_2(Q_2^2-Q_4)$ and $6p_3(2Q_4-Q_2^2)$. Subtracting these from $(\sum_a y_a)^4=0$ with all $p_j$ set to one gives the distinct-index term $p_4(3Q_2^2-6Q_4)$. Since $p_j\le p_1\le2a_s$, the fourth moment is at most $Ca_s\|y\|^4$. Cauchy--Schwarz yields $\mathbb E|L|^3\le Ca_s\|y\|^3$. Taylor's inequality and the logarithm series therefore give

\begin{equation*}
\log f_s(y)=-\frac{a_s}{2}\|y\|^2+O(a_s\|y\|^3).
\end{equation*}

The branch passes through zero at $y=0$. The extra series error $O(a_s^2\|y\|^4)$ is absorbed since $a_s\le1$ and $\|y\|\le R$. Thus the centered integrand is $\exp(-\Lambda\|y\|^2/2+B(y))$, with

\begin{equation*}
\Lambda=\sum_j a_{s_j},\qquad |B(y)|\le C\Lambda\|y\|^3,\qquad
(n-e)/d\le\Lambda\le n.
\end{equation*}

Use $|e^B-1|\le |B|e^{|B|}$ and $\|y\|\le1/d$. The relative error integrated over the ball, normalized by the full target Gaussian integral, is at most

\begin{equation*}
C\Lambda(1-2C/d)^{-(d-1)/2}\mathbb E\|Y\|^3
=O(d^{3/2}/\sqrt\Lambda)=O(d^2/\sqrt n)=o(1),
\end{equation*}

where $Y$ is Gaussian on $V$ of precision $\Lambda(1-2C/d)$. The normalization ratio is bounded and the moment bound follows from the Gaussian fourth moment. The target Gaussian mass outside the ball is $O(d/(\Lambda R^2))=O(d^4/n)=o(1)$ by Markov's inequality. Its full torus-density-normalized integral is $\sqrt d(2\pi\Lambda)^{-(d-1)/2}\ge e^{-C_1d\log n}$ for fixed $C_1$. As $n/\log n$ increases for $n>e$, the ratio $n/(d^4\log n)$ has lower limit at least $C_*/4$ in the stated range. Choose $C_*$ large enough that $c'n/d^3\ge 2C_1 d\log n$ there for all large $d$. The exterior is then negligible relative to the Gaussian. All estimates are uniform in the profile and in $r$.

The profile dependence in the determinant can now be discarded, since

\begin{equation*}
0\le \Lambda-\Lambda_*(i)
=\frac{\sum_j s_j(r-s_j)}{d(d-1)}\le\frac{re}{d(d-1)},
\qquad
d\log(\Lambda/\Lambda_*(i))
\le\frac{dre}{di(d-r)}
=O(d^3/n).
\end{equation*}

In the last estimate use $di=nr-e\sim nr$ and $d-r\ge1$. Averaging over profiles in $A_i$ with weights proportional to $\prod_j\binom d{s_j}$ gives the asserted count, including $e=0$. The profile event describes exactly the balanced arrays conditional on the row sizes. The logarithmic first differences of the determinant factor at these quotas are $O(d/i)=O(d^2/(nr))=o(1)$, so the same is true of its second differences. Together with the uniform relative error, this means all required fixed-offset logarithmic differences may be taken on $A_i$.

\subsection{Scalar quota differences}

Put $b=d^2/n$. At a real $x\in(0,T)$, let the size law with positive tilt $w_x$ have mean $dx/n$. The mean is strictly increasing in $\log w_x$ since its derivative is the positive variance $V_x$. As $w_x$ spans the positive axis the mean spans $(0,r)$. The deficit $J=r-S$ has mean $h_x=d(T-x)/n$ and weights $B_j z^j$ for $B_j=\binom d{r-j}$ and $z=1/w_x$. Write $\ell=B_1z/B_0$. The forward weight ratios decrease in $j$, so $B_j z^j/B_0\le\ell^j$. If $h_x=o(1)$, the mass at zero is at least $1-h_x$ while the mass at one is at most $h_x$, so $\ell\le h_x/(1-h_x)$. Summing the resulting geometric majorant and retaining the first two masses gives, uniformly in the cap,

\begin{equation*}
h_x=\ell+O(\ell^2),\qquad
V_x=h_x+O(h_x^2),\qquad
\mathbb E|J-h_x|^3=O(h_x),\qquad
\mathbb P(J=0)\mathbb P(J=1)\asymp h_x.
\end{equation*}

Here all raw moments up to order three from $J\ge2$ are $O(\ell^2)$, since $\sum_{j\ge2}j^m\ell^j=O(\ell^2)$ for $\ell\le1/2$. The derivative of $V_x$ with respect to the size mean is the signed third centered size moment divided by variance, hence is uniformly $O(1)$ in this regime.

Suppose $(D-k)/d\to c\in(0,\infty)$. On any bounded real window around $k$, $h_x\sim cb$ and $nh_x\sim cd^2$. At integer $i$ in that window, the sum of $n$ tilted deficits has integral mean $e=nr-di$. The two-mass product estimate implies that the modulus of the deficit characteristic function is at most $\exp(-c_2h_i t^2)$ for $|t|\le\pi$ by comparing two independent copies and retaining their differences $\pm1$. Rescale Fourier inversion at the mean by $t=v/\sqrt{nV_i}$. The centered expansion of the logarithm on bounded $v$ gives $-v^2/2+O(|v|^3/\sqrt{nh_i}+v^4/n)$ after multiplication by $n$. The modulus bound supplies an integrable Gaussian majorant independent of the parameters. Thus the central probability is $(2\pi nV_i)^{-1/2}(1+o(1))$ at each fixed integer shift.

For $E(x)=n\log H(w_x)-dx\log w_x$, one has $E'=-d\log w_x$ and $E''=-b/V_x\to-1/c$ uniformly on the window. Also $E'''=O(b(d/n)/h_x^2)=O(1/d)$ by the derivative bound just proved. The tilt identity yields $\log A_i=E(i)-\tfrac12\log(2\pi nV_i)+o(1)$, and neighboring $V_i$ have ratio tending to one. The central second difference of $E$ gives $\log\Gamma_{k+j}\to1/c$, and $(E(k-1)-E(k+1))/2=-E'(k)+O(1/d)$ gives the claimed logarithm of $t_k$. The determinant factor has already been accounted for.

\subsection{Distances sublinear in dimension}

Use scalar deficit counts $U_e=[z^e](\sum_j B_j z^j)^n$. They are log-concave relative to $\binom{n+e-1}e$. To see this, let $C_m$ be the analogous totals for $n-1$ rows, extended by zero. They are log-concave with interval support by convolution. For increasing $e$, the law on $j$ with weights $(j+1)B_j C_{e-j}$ increases in likelihood ratio order: the common-support ratio is proportional to $C_{e+1-j}/C_{e-j}$, increasing with $j$, and the support endpoints both increase weakly. Under this law the expectation of the decreasing sequence $B_{j+1}/B_j$, extended by zero at $r$, equals $(e+1)U_{e+1}/[(n+e)U_e]$. The numerator comes from differentiating the row polynomial power, and the denominator from exchangeability of the rows conditional on their sum. Thus the normalized forward ratio decreases, proving the assertion.

At an internal quota with $(D-k)/d\to0$, $e=d(T-k)\ge d$ and $e=o(d^2)$. The log curvature of $A_k=U_e$ at quota step one is at least the log curvature of its negative-binomial normalizer at deficit step $d$. Writing the binomial forward logarithmic increment as $g(v)=\log[(n+v-1)/v]$, this gives

\begin{equation*}
\log\frac{A_k^2}{A_{k-1}A_{k+1}}
\ge\sum_{j=1}^d \big(g(e-d+j)-g(e+j)\big)
\ge\frac{d^2(n-1)}{(e+d)(n+e+d)}\longrightarrow\infty .
\end{equation*}

Indeed $-g'(v)=(n-1)/[v(n+v-1)]$ and each difference integrates over length $d$. This argument includes the case $e=d$, where its upper-quotient neighbor has exact saturation. The determinant comparison transfers the divergence to $\log\Gamma_k$.
